% 解説プリントのコメント（saiten が書く。./tools/kaisetsu.ps1 がこのファイルからプリントを作る）
% 熱物理学 第8章 問題1・2・3・5・6・7（2026-11-16 解答、スキャン 3枚）
%% matome
\begin{daiji}
  \item 偏微分は「何を固定するか」までが定義である。$C_{V}=(\partial U/\partial T)_{V}$ であり、$(\partial U/\partial T)_{p}$ と値が一致するのは $U$ が $T$ だけの関数のときに限る（問題6 (4)）。
  \item 「基準 $(T_{0},p_{0})$ での値を $S_{0}$ とする」と指定されたら、積分の下端を基準状態にして $\ln(T/T_{0})$ の形で書く。対数の中は無次元にする（問題6 (3)）。
  \item 1つの小問に問いが2つあるときは後半まで答える（問題5 (3) の $C_{V}$）。「有限に留まることを述べよ」には、有限になる理由（$C_{V}/T$ が $T\to0$ で積分できる）まで書く（問題2 (3)）。
\end{daiji}
\begin{maemo}
  \item 小問の後半の答え落とし：第3章 問題2 (3)（2026-11-12）で原始関数が未着手、問題4 (5) で理想気体との比較が無かった。今回も問題5 (3) で $C_{V}$ が無い。解き終えたら問題文の「求めよ」「表せ」「確かめよ」を数え、それぞれに答えがあるか印を付ける。
  \item 固定する変数の取り違え：第3章 問題6 (1)（2026-11-12）で「固定した変数を考えていない」とした。今回は $(\partial U/\partial T)_{p}=C_{V}$ と書いた。偏微分を書くたびに、添字が問題の定義（$C_{V}$ なら $V$ 一定）と合っているかを確かめる。
\end{maemo}
%% mondai 2
\begin{kotaeran}
  \machigai{3}{○}{記述} $T_{\mathrm{m}}\to 3T_{2}/4$ と、理想気体では $\ln(T_{2}/T_{1})$ が $T_{1}\to0$ で発散することは正しい。肝心の「$\Delta S$ が有限に留まる」ことと、その理由が書かれていない。(1) から $\Delta S\to AT_{2}^{3}/3$（有限）であり、それは被積分関数 $C_{V}/T=AT^{2}$ が $T\to0$ で $0$ になって積分できるから（比熱が $T\to0$ で十分速く $0$ になる）である。
\end{kotaeran}
\begin{rikai}
  \item $\Delta S=\int C_{V}/T\,\dd T$ の下端 $T_{1}\to0$ での振る舞いは、被積分関数で決まる。$C_{V}=AT^{3}$ なら $AT^{2}$ で有限、比熱一定の理想気体なら $nc_{V}/T$ で、$\ln T_{1}\to-\infty$ と発散する。2つを並べて「被積分関数の違い」として述べる。
\end{rikai}
%% mondai 5
\begin{kotaeran}
  \machigai{3}{△}{未着手} $(\partial U/\partial V)_{T}=\frac{1}{3}(AT/4)^{4/3}(n/V)^{2/3}$ と $(\partial p/\partial T)_{V}$ を求め、エネルギー方程式が成り立つことを確かめたところ（両辺とも $p/3$）は正しい。問われている定積比熱 $C_{V}$ が書かれていない。(2) の $U(T,V)$ を $V$ 一定で $T$ について微分すれば $C_{V}=(4/3)(A/4)^{4/3}n^{2/3}V^{1/3}T^{1/3}$ である。
\end{kotaeran}
\begin{rikai}
  \item $C_{V}=(\partial U/\partial T)_{V}$ は、$U$ を $(T,V)$ の関数として書けば微分1回で出る。(2) で $U(T,V)$ を求めさせたのはこのためである。
  \item 検算：$U\propto T^{4/3}$ なので $C_{V}=4U/(3T)>0$ になっているかを見る。
\end{rikai}
%% mondai 6
\begin{kotaeran}
  \machigai{3}{○}{立式} $\dd S=nc_{p}\,\dd T/T-nR\,\dd p/p$ は正しい。積分した結果を $S(T,p)=nc_{p}\ln T-nR\ln p+S_{0}$ と書いたが、問題は $S_{0}$ を基準 $(T_{0},p_{0})$ でのエントロピーと指定している。この式に $T=T_{0}$、$p=p_{0}$ を入れても $S_{0}$ にならず、$S_{0}$ がただの積分定数になっている。また $\ln T$、$\ln p$ は次元のある量の対数である。下端を $(T_{0},p_{0})$ にとって $S=nc_{p}\ln(T/T_{0})-nR\ln(p/p_{0})+S_{0}$ と書く。
  \machigai{4}{○}{記述} $U=H-pV=n(c_{p}-R)T+H_{0}$ は正しいが、続けて $(\partial U/\partial T)_{p}=n(c_{p}-R)=C_{V}$ と書いた。$C_{V}$ は $(\partial U/\partial T)_{V}$ であり、圧力一定の微分ではない。ここで値が一致するのは $U$ が $T$ だけの関数だからで、「$U$ は $T$ だけの関数なので $(\partial U/\partial T)_{V}=\dv{U}{T}=n(c_{p}-R)$」と書く。
\end{kotaeran}
\begin{rikai}
  \item 2つの微分の関係は $(\partial U/\partial T)_{p}=(\partial U/\partial T)_{V}+(\partial U/\partial V)_{T}(\partial V/\partial T)_{p}$ である。一般には第2項があるので $(\partial U/\partial T)_{p}\neq C_{V}$。この気体では $U$ が $V$ によらず $(\partial U/\partial V)_{T}=0$ なので一致する。比熱の定義は $C_{V}=(\partial U/\partial T)_{V}$、$C_{p}=(\partial H/\partial T)_{p}$ と、固定する変数まで覚える。
  \item 対数の中は無次元にする。基準状態で割った $\ln(T/T_{0})$ の形にすれば、$T=T_{0}$、$p=p_{0}$ を代入して $S=S_{0}$ になることで検算できる。
\end{rikai}
%% mondai 7
\begin{kotaeran}
  \machigai{3}{○}{記述} $s^{**}=s^{*}+c_{V}\ln(T^{**}/T^{*})+R\ln(\rho^{*}/\rho^{**})$ と、それが元の式での $(T^{**},\rho^{**})$ における単位物質量当たりのエントロピーであることは正しい。最後の2状態の差を $nc_{V}\ln(T_{2}/T_{1})+nR\,V_{2}/V_{1}$ と書き、第2項の $\ln$ が落ちている。正しくは $nR\ln(V_{2}/V_{1})$。
\end{kotaeran}
\begin{rikai}
  \item 理想気体のエントロピーの差は、各項が「定数 $\times$ $\ln$（無次元の比）」の形である。書き終えたら同じ状態（$T_{2}=T_{1}$、$V_{2}=V_{1}$）を入れて差が $0$ になるかを確かめる。$nR\,V_{2}/V_{1}$ のままだと $nR$ が残ってしまう。
\end{rikai}
